% Part: many-valued-logic
% Chapter: three-valued-logics
% Section: lukasiewicz
%READERNOTE{SOL6-A140: दिलेल्या मोडल कोष्टकांत मधल्या U मूल्यावरील विरोधाभासाची संभाव्यता T होते; तिचे निषेधन F होते, मूळात दिलेले U नाही. मॅट्रिक्सचा चार-संयोजकीय भाषाखंड स्पष्ट केला. मूळ 1920 चा व्याख्यानवृत्तांत प्रत्यक्ष पाहून ऐतिहासिक वर्ष दुरुस्त केले; उद्धरणातील noon म्हणजे दुपारी बारा वाजता.}

\documentclass[../../../include/open-logic-section]{subfiles}

\begin{document}

\olfileid{mvl}{thr}{luk}

\olsection{\L ukasiewicz यांचे तर्कशास्त्र}

बहुमूल्य तर्कशास्त्राचे प्रकाशित झालेले आणि सविस्तर मांडलेले
सर्वांत आरंभीच्या प्रस्तावांपैकी एक पोलिश तत्त्वज्ञ Jan \L ukasiewicz
यांनी 1920 मध्ये मांडला.\footnote{Jan \L ukasiewicz,
``O logice trójwartościowej, \emph{Ruch Filozoficzny} 5
(1920), पृ. 170--171; 19 जून 1920 च्या व्याख्यानाचा वृत्तांत.
मूळ अंकाचा स्कॅन: \url{https://kpbc.umk.pl/Content/247342/PDF/EOD_OPEN_006_12_nr_09.pdf}.} अरिस्टॉटलच्या सागरी युद्धाच्या समस्येतून
\L ukasiewicz यांना प्रेरणा मिळाली: \emph{आज} ``उद्या सागरी युद्ध
होईल'' हे विधान सत्यही नाही आणि असत्यही नाही असे दिसते; त्याचे
सत्यतामूल्य अजून निश्चित झालेले नाही. अशा ``भविष्यातील
परिस्थितीवश'' विधानांसाठी \L ukasiewicz यांनी तिसरे सत्यतामूल्य
आणण्याचा प्रस्ताव मांडला.
\begin{quote}
पुढील वर्षी एखाद्या ठराविक क्षणी, उदाहरणार्थ 21 डिसेंबरला दुपारी बारा वाजता,
मी वॉर्सामध्ये असेन, हे सध्या होकारार्थी किंवा नकारार्थी अशा
कोणत्याही प्रकारे निश्चित झालेले नाही, असे मी विरोधाभासाशिवाय
गृहीत धरू शकतो. म्हणून त्या वेळी मी वॉर्सामध्ये असेन हे संभाव्य
आहे, पण अनिवार्य नाही. या गृहितकावर ``पुढील वर्षी 21 डिसेंबरला
दुपारी बारा वाजता मी वॉर्सामध्ये असेन'' हे विधान सध्या सत्यही नाही आणि
असत्यही नाही. कारण ते आताच सत्य असेल, तर भविष्यात मी
वॉर्सामध्ये असणे अनिवार्य ठरेल; हे गृहितकाशी विसंगत आहे.
दुसरीकडे ते आताच असत्य असेल, तर भविष्यात मी वॉर्सामध्ये असणे
अशक्य ठरेल; हेही गृहितकाशी विसंगत आहे. म्हणून विचाराधीन
विधान या क्षणी सत्यही नाही आणि असत्यही नाही; त्याला ``0''
म्हणजे असत्यता आणि ``1'' म्हणजे सत्यता यांपेक्षा वेगळे तिसरे
मूल्य असले पाहिजे. या मूल्याला आपण ``$\frac{1}{2}$.'' असे
दर्शवू शकतो. ते ``संभाव्य'' याचे प्रतिनिधित्व करते आणि
``सत्य'' व ``असत्य'' यांच्याबरोबर तिसरे मूल्य म्हणून सामील होते.
\end{quote}
\L ukasiewicz यांच्या तिसऱ्या सत्यतामूल्यासाठी आपण $\Undef$
वापरू.\footnote{\L ukasiewicz येथे ``संभाव्य'' हा शब्द आज रूढ
नसलेल्या अर्थाने, म्हणजे संभाव्य पण अनिवार्य नाही, या अर्थाने
वापरतात.}

या अर्थानुसार $\lnot$, $\land$ आणि $\lor$ या संयोजकांची
सत्यता-फलने ठरवणे सोपे आहे: भविष्यातील परिस्थितीवश विधानाचे
निषेधनही भविष्यातील परिस्थितीवश विधान असते; म्हणून
$\tf{\lnot}(\Undef) = \Undef$. संधीमधील एक घटक अनिर्धारित
असेल आणि दुसरा सत्य असेल, तर ती संधीही अनिर्धारित असते.
कारण परिस्थितीवश घटक पुढे कसा ठरतो त्यानुसार संधी सत्यही
ठरू शकते आणि असत्यही ठरू शकते. म्हणून \[
    \tf{\land}(\True, \Undef) =
\tf{\land}(\Undef, \True) =
\Undef.
\]
पण दुसरा घटक असत्य असेल, तर संधी सत्य ठरू शकत नाही;
म्हणून \[\tf{\land}(\False, \Undef) =
\tf{\land}(\Undef, \False) = \False.\]
दोन्ही निविष्टे $\True$ किंवा $\False$ ही निश्चित सत्यतामूल्ये
असतील, तर उरलेली फलनमूल्ये अभिजात तर्कशास्त्राप्रमाणेच असतात.

सशर्त विधानाच्या बाबतीत परिस्थिती थोडी अधिक गुंतागुंतीची आहे.
$q$ हे भविष्यातील परिस्थितीवश विधान आहे असे समजा. $p$ असत्य
असेल, तर $q$ पुढे कसेही ठरले तरी $p \lif q$ सत्य असेल; म्हणून
$\tf{\lif}(\False, \Undef) = \True$ असे ठेवावे. तसेच $p$
सत्य असेल, तर $q$ पुढे कसेही ठरले तरी $q \lif p$ सत्य असेल;
म्हणून $\tf{\lif}(\Undef, \True) = \True$. $p$ सत्य असेल,
तर $p \lif q$ सत्य किंवा असत्य ठरू शकते; म्हणून
$\tf{\lif}(\True, \Undef) = \Undef$. त्याचप्रमाणे $p$
असत्य असेल, तर $q \lif p$ सत्य किंवा असत्य ठरू शकते;
म्हणून $\tf{\lif}(\Undef, \False) = \Undef$. आता $p$
आणि $q$ दोन्ही भविष्यातील परिस्थितीवश विधाने असण्याचे
प्रकरण उरते. मूळ प्रेरणेनुसार या प्रकरणात खरे तर $\Undef$
हेच मूल्य द्यायला हवे. मात्र तसे केल्यास $!A \lif !A$
हे \emph{सर्वतः सत्य सूत्र राहणार नाही}. $!A \lor \lnot !A$
आणि $\lnot(!A \land \lnot !A)$ यांचा त्याग करायला
\L ukasiewicz यांना अडचण वाटली नाही; पण $!A \lif !A$
चा त्याग करणे त्यांना मान्य नव्हते. म्हणून त्यांनी
$\tf{\lif}(\Undef, \Undef) = \True$ असे ठरवले.

\begin{defn}\ollabel{def:lukasiewicz}
त्रिमूल्य \L ukasiewicz तर्कशास्त्र पुढील मॅट्रिक्सने परिभाषित होते:
\begin{enumerate}
  \item मानक विधानीय भाषा $\Lang L_0$ चा केवळ
  $\lnot$, $\land$, $\lor$, $\lif$ हे संयोजक असलेला खंड.
  \item सत्यतामूल्यांचा संच $V = \{\True, \Undef, \False\}$.
  \item $\True$ हेच एकमेव निर्दिष्ट मूल्य आहे, म्हणजे $V^+ = \{\True\}$.
  \item सत्यता-फलने पुढील कोष्टकांनी दिली आहेत:
  \begin{center}
    \begin{tabular}{c|c}
      $\tf{\lnot}$ & \\
      \hline
      $\True$ & $\False$ \\
      $\Undef$ & $\Undef$ \\
      $\False$ & $\True$
    \end{tabular}
    \quad
    \begin{tabular}{c|ccc}
      $\tf{\land}[\LogLuk[3]]$ & $\True$ & $\Undef$ & $\False$ \\
      \hline
      $\True$ & $\True$ & $\Undef$ & $\False$ \\
      $\Undef$ & $\Undef$ & $\Undef$ & $\False$\\
      $\False$ & $\False$ & $\False$ & $\False$
    \end{tabular}
    \\[2ex]
    \begin{tabular}{c|ccc}
      $\tf{\lor}[\LogLuk[3]]$ & $\True$ & $\Undef$ & $\False$ \\
      \hline
      $\True$ & $\True$ & $\True$ & $\True$ \\
      $\Undef$ & $\True$ & $\Undef$ & $\Undef$ \\
      $\False$ & $\True$ & $\Undef$ & $\False$
    \end{tabular}
    \quad
    \begin{tabular}{c|ccc}
      $\tf{\lif}[\LogLuk[3]]$ & $\True$ & $\Undef$ & $\False$ \\
      \hline
      $\True$ & $\True$ & $\Undef$ & $\False$ \\
      $\Undef$ & $\True$ & $\True$ & $\Undef$  \\
      $\False$ & $\True$ & $\True$ & $\True$
    \end{tabular}
  \end{center}
\end{enumerate}
\end{defn}

सहज लक्षात येते की केवळ $\lnot$, $\land$ आणि $\lor$
असलेल्या कोणत्याही !!{formula}~$!A$ मधील सर्व
!!{propositional variable}s ना $\Undef$ नेमले, तर त्या
सूत्राचे सत्यतामूल्य~$\Undef$ असेल. म्हणून, उदाहरणार्थ,
अभिजात तर्कशास्त्रातील $p \lor \lnot p$ आणि $\lnot(p
\land \lnot p)$ ही सर्वतः सत्य सूत्रे $\LogLuk[3]$
मध्ये सर्वतः सत्य नाहीत; कारण $\pAssign v(p) = \Undef$
असताना $\pValue{v}(!A) = \Undef$.

संबंधित प्रत्येक विधानीय चलासाठी $\pAssign v(p) = \True$ किंवा
$\False$ असलेल्या !!{valuation}s मध्ये $\pValue v(!A)$
हे त्या सूत्राच्या अभिजात सत्यतामूल्याशी जुळते.

\begin{prop}
  $!A$ मधील प्रत्येक $p$ साठी $\pAssign v(p) \in \{\True, \False\}$ असेल, तर
  $\pValue v(!A)[\LogLuk[3]] = \pValue v(!A)[\LogCL]$.
\end{prop}

\begin{prob}\label{mvl:thr:luk:prob:luk-iff} असे समजा की आपण
  $\pValue v(!A \liff !B) = \pValue v((!A \lif !B) \land (!B \lif
  !A))$ अशी $\LogLuk[3]$ मध्ये व्याख्या करतो. मग $\liff$
  चे सत्यताकोष्टक कसे असेल?
\end{prob}

अभिजात तर्कशास्त्रातील अनेक सर्वतः सत्य सूत्रे \LogLuk[3]
मध्येही \emph{सर्वतः सत्य आहेत}; उदाहरणार्थ,
$\lnot p \lif (p \lif q)$. अभिजात तर्कशास्त्राप्रमाणेच
हे सत्यताकोष्टकाने पडताळता येते:
\begin{center}
\begin{tabular}{cc|cccccc}
  $p$ & $q$ & $\lnot$ & $p$ & $\lif$ & $\smash{(}p$ & $\lif$ & $q\smash{)}$ \\ \hline
  \True & \True & \False & \True & \True & \True & \True & \True \\
  \True & \Undef & \False & \True & \True & \True & \Undef & \Undef \\
  \True & \False & \False & \True & \True & \True & \False & \False \\
  \Undef & \True & \Undef & \Undef & \True & \Undef & \True & \True \\
  \Undef & \Undef & \Undef & \Undef & \True & \Undef & \True & \Undef \\
  \Undef & \False & \Undef & \Undef & \True & \Undef & \Undef & \False \\
  \False & \True & \True & \False & \True & \False & \True & \True \\
  \False & \Undef & \True & \False & \True & \False & \True & \Undef \\
  \False & \False & \True & \False & \True & \False & \True & \False \\
\end{tabular}
\end{center}

\begin{prob}
  पुढील सूत्रे \LogLuk[3] मध्ये सर्वतः सत्य आहेत हे दाखवा:
  \begin{enumerate}
    \item $p \lif (q \lif p)$
    \item\label{mvl:thr:luk:prob:luk-taut-2} $\lnot(p \land q) \liff (\lnot p \lor \lnot q)$
    \item\label{mvl:thr:luk:prob:luk-taut-3} $\lnot(p \lor q) \liff (\lnot p \land \lnot q)$
  \end{enumerate}
  (\olref[mvl][thr][luk]{prob:luk-taut-2} आणि
  \olref[mvl][thr][luk]{prob:luk-taut-3} मध्ये $!A \liff !B$
  हा $(!A \lif !B) \land (!B \lif !A)$ चा संक्षेप माना किंवा
  \olref[mvl][thr][luk]{prob:luk-iff} साठीचे तुमचे उत्तर वापरा.)
\end{prob}

\begin{prob}
  पुढील अभिजात सर्वतः सत्य सूत्रे \LogLuk[3] मध्ये सर्वतः सत्य
  नाहीत हे दाखवा:
  \begin{enumerate}
    \item $(\lnot p \land p) \lif q$
    \item $((p \lif q) \lif p) \lif p$
    \item $(p \lif (p \lif q)) \lif (p \lif q)$
  \end{enumerate}
\end{prob}

त्यामुळे अभिजात तर्कशास्त्रातील सर्व सर्वतः सत्य सूत्रे
$\LogLuk[3]$ मध्ये सर्वतः सत्य नसली, तरी प्रत्येक
!!{valuation} मध्ये त्यांना निदान $\True$ किंवा $\Undef$
यांपैकी एक मूल्य मिळेल, असे वाटू शकते. पण तसे नाही.
पुढील सूत्र हे प्रतिउदाहरण आहे:
\[
  \lnot(p \lif \lnot p) \lor \lnot(\lnot p \lif p)
\]
$p$ चे मूल्य~$\Undef$ असेल, तर या सूत्राचे मूल्य $\False$ येते.

\begin{prob}
  पुढीलपैकी कोणते संबंध \L ukasiewicz तर्कशास्त्रात खरे ठरतात?
  प्रत्येकासाठी सत्यताकोष्टक द्या.
  \begin{enumerate}
    \item $p, p \lif q \Entails q$
    \item $\lnot\lnot p \Entails p$
    \item $p \land q \Entails p$
    \item $p \Entails p \land p$
    \item $p \Entails p \lor q$
  \end{enumerate}
\end{prob}

\L ukasiewicz यांना आपल्या त्रिमूल्य पद्धतीच्या आधारे
संभाव्यतेचे तर्कशास्त्र उभारायचे होते. त्यासाठी त्यांनी ``$!A$
संभाव्य आहे'' याकरिता एकस्थानी संयोजक $\Diamond !A$ आणि
``$!A$ अनिवार्य आहे'' याकरिता त्याला अनुरूप $\Box !A$ घेतले:
\begin{center}
  \begin{tabular}{c|c}
    $\tf{\Diamond}$ & \\
    \hline
    $\True$ & $\True$ \\
    $\Undef$ & $\True$ \\
    $\False$ & $\False$
  \end{tabular}
  \quad
  \begin{tabular}{c|c}
    $\tf{\Box}$ & \\
    \hline
    $\True$ & $\True$ \\
    $\Undef$ & $\False$ \\
    $\False$ & $\False$
  \end{tabular}
\end{center}
म्हणजे $p$ संभाव्य आहे, जर आणि फक्त जर ते असत्य असल्याचे
आधीच निश्चित झालेले नसेल; आणि $p$ अनिवार्य आहे, जर आणि
फक्त जर ते सत्य असल्याचे आधीच निश्चित झालेले असेल.

\begin{prob}
  $\Box$ आणि~$\Diamond$ यांची सत्यताकोष्टके जोडून विस्तारलेल्या
  \LogLuk[3] मध्ये $\Box p \liff \lnot\Diamond \lnot p$ आणि
  $\Diamond p \liff \lnot \Box \lnot p$ ही सर्वतः सत्य सूत्रे
  आहेत हे दाखवा.
\end{prob}

पण या प्रस्तावित मोडल तर्कशास्त्राच्या मर्यादा लवकरच स्पष्ट
झाल्या: पुढे काहीही घडले तरी $p \land \lnot p$ कधीच सत्य
ठरू शकत नाही. म्हणून सध्या त्याचे मूल्य निश्चित नसले
(आणि त्यामुळे ते अनिर्धारित असले), तरी ते अशक्य मानायला
हवे; म्हणजे $\lnot \Diamond(p \land \lnot p)$ हे सर्वतः सत्य
सूत्र असायला हवे. तथापि, $\pAssign v(p) = \Undef$ असेल,
तर $\pValue v(\lnot \Diamond(p \land \lnot p)) = \False$.
कारण या वेळी $p \land \lnot p$ चे मूल्य $\Undef$,
त्याच्या $\Diamond$ चे मूल्य $\True$, आणि शेवटच्या निषेधनाचे
मूल्य $\False$ असते.
संभाव्यता आणि अनिवार्यता यांसारखे मोडल भेद व्यक्त करण्यासाठी
दोन सत्यतामूल्ये पुरेशी नाहीत, हा \L ukasiewicz यांचा निष्कर्ष
बरोबर होता; पण तिसरे सत्यतामूल्य जोडणेही पुरेसे नाही.

\end{document}
